\DocumentMetadata{
  pdfversion=2.0,
  pdfstandard=ua-2,
  lang=en-US,
  tagging=on,
  tagging-setup={math/setup=mathml-SE,tabsorder=structure}
}
\documentclass[11pt,letterpaper]{article}
\usepackage[margin=0.68in,includefoot]{geometry}
\usepackage{newcomputermodern}
\usepackage{amsmath,amscd,mathtools}
\usepackage{tikz,adjustbox}
\providecommand{\square}{\mdlgwhtsquare}
\usepackage[hidelinks]{hyperref}
\hypersetup{
  pdftitle={Algebra PhD Exam, May 2020},
  pdfauthor={University of Florida Department of Mathematics},
  pdfsubject={Graduate mathematics examination},
  pdfdisplaydoctitle=true,
  bookmarksopen=true
}
\setcounter{secnumdepth}{0}
\setlength{\parindent}{0pt}
\setlength{\parskip}{0.28em}
\setlength{\emergencystretch}{3em}
\setlength{\leftmargini}{2.15em}
\setlength{\leftmarginii}{2.65em}
\setlength{\leftmarginiii}{3em}
\pagestyle{empty}
\raggedbottom
\allowdisplaybreaks
\NewTaggingSocketPlug{block/list/label}{examproblem}{%
  \tagstructbegin{tag=Lbl}\tagstructend%
  \tagstructbegin{tag=\LBody}%
  \tagstructbegin{tag=\ProblemHeadingTag,title-o={\ProblemHeadingText}}%
  \tagmcbegin{tag=\ProblemHeadingTag,actualtext={\ProblemHeadingText}}%
  #2%
  \tagmcend%
  \tagstructend%
}
\newenvironment{examproblems}[2]{%
  \def\ProblemHeadingTag{#2}%
  \AssignTaggingSocketPlug{block/list/label}{examproblem}%
  \begin{enumerate}%
  \setcounter{enumi}{#1}%
  \renewcommand{\labelenumi}{\textbf{\theenumi.}}%
  \setlength{\topsep}{0.35em}%
  \setlength{\partopsep}{0pt}%
  \setlength{\itemsep}{0.48em}%
  \setlength{\parsep}{0.18em}%
}{\end{enumerate}}
\newenvironment{examsubparts}{%
  \AssignTaggingSocketPlug{block/list/label}{default}%
  \begin{enumerate}%
  \setlength{\topsep}{0.18em}%
  \setlength{\partopsep}{0pt}%
  \setlength{\itemsep}{0.16em}%
  \setlength{\parsep}{0.1em}%
}{\end{enumerate}}
\newenvironment{examinstructions}{%
  \par\begingroup\itshape
}{%
  \par\endgroup\vspace{0.18em}
}
\makeatletter
\renewcommand\subsection{\@startsection{subsection}{2}{0pt}%
  {-1.25ex plus -.25ex minus -.1ex}%
  {0.55ex plus .1ex}%
  {\normalfont\large\bfseries}}
\renewcommand{\@maketitle}{%
  \newpage\null\vskip -1em%
  {\footnotesize\bfseries
    \href{https://www.ufl.edu/}{University of Florida}, \href{https://math.ufl.edu/}{Department of Mathematics}\par}%
  \vskip -0.5em%
  \tagstructbegin{tag=H1,title={Algebra PhD Exam, May 2020}}%
  \tagpdfparaOff%
  \tagmcbegin{tag=H1}%
    {\LARGE\bfseries\href{https://gma.math.ufl.edu/exams/algebra-phd/}{Algebra PhD Exam}, May 2020\par}%
  \tagmcend%
  \tagpdfparaOn%
  \tagstructend%
  \vskip 0.25em%
  \tagmcbegin{artifact}%
    \hrule height 1pt%
  \tagmcend%
  \vskip 1em%
}
\makeatother
\title{Algebra PhD Exam, May 2020}
\author{}
\date{}
\begin{document}
\maketitle
\thispagestyle{empty}
\begin{examinstructions}
Answer seven problems. (If you turn in more, the first seven will be graded.) Within reason, you may quote theorems as long as you state them clearly.

\par
Below, ring means associative ring with identity, and module means unital module unless otherwise specified.
\end{examinstructions}
\begin{examproblems}{0}{H2}
\def\ProblemHeadingText{Problem 1.}
\item
Let \(q>1\) be a power of a prime~\(p\). Prove that there exists some field~\(F\) with \(|F|=q\).
\def\ProblemHeadingText{Problem 2.}
\item
Let~\(F\) be the splitting field of \(p(x)=x^3+5x+5\) over~\(\mathbb{Q}\). Prove that there is a unique intermediate field~\(K\) such that \(K\ne\mathbb{Q}\), \(K\ne F\), and \(K/\mathbb{Q}\) is Galois.
\def\ProblemHeadingText{Problem 3.}
\item
Let~\(A\) be a commutative ring with identity. Suppose that~\(P\) and~\(Q\) are projective~\(A\)-modules. Prove that \(P\otimes_A Q\) is a projective~\(A\)-module.
\def\ProblemHeadingText{Problem 4.}
\item
Let~\(\mathcal{C}\) be a concrete category. Define what is meant by a free object in~\(\mathcal{C}\). Let~\(\mathcal{G}\) be the category of all finite groups. Prove that the free objects in~\(\mathcal{G}\) are exactly the trivial groups.
\def\ProblemHeadingText{Problem 5.}
\item
Let~\(R\) be a ring. Prove that the following two conditions are equivalent for a left~\(R\)-module~\(P\). (You may not assume any properties of projective modules.)
\begin{examsubparts}
\item[\textnormal{(a)}]
Every short exact sequence of~\(R\)-modules
\[
0\to A\to B\to P\to 0
\]
 splits;
\item[\textnormal{(b)}]
There is a free~\(R\)-module~\(F\) and an~\(R\)-module~\(K\) such that \(F\simeq K\oplus P\).
\end{examsubparts}
\def\ProblemHeadingText{Problem 6.}
\item
Give an example of a polynomial \(f(X)\in\mathbb{Q}[X]\) whose Galois group is isomorphic to \(S_5\). Prove that this is the case.
\def\ProblemHeadingText{Problem 7.}
\item
Prove: if a Dedekind domain has only a finite number of nonzero prime ideals then it is a principal ideal domain. (Hint: prove first that each prime ideal is principal, use the Chinese Remainder Theorem.)
\def\ProblemHeadingText{Problem 8.}
\item
Let~\(R\) be an integral domain with field of fractions~\(K\). For \(a\in K\), let \(D(a)=\{r\in R\mid ra\in R\}\).
\begin{examsubparts}
\item[\textnormal{(a)}]
Show that \(D(a)\) is an ideal and that \(D(a)=R\) if and only if \(a\in R\).
\item[\textnormal{(b)}]
For each prime ideal~\(P\), let \(R_P\) denote the elements of~\(K\) which can be represented as a fraction having denominator not in~\(P\). Show that \(a\in R_P\) if and only if \(D(a)\nsubseteq P\).
\item[\textnormal{(c)}]
Deduce that \(\bigcap_P R_P=R\), where the intersection is taken over the maximal ideals of~\(R\).
\end{examsubparts}
\def\ProblemHeadingText{Problem 9.}
\item
For this question you may not quote Nakayama’s Lemma. Let~\(R\) be a commutative local ring with~\(1\) and~\(J\) its unique maximal ideal.
\begin{examsubparts}
\item[\textnormal{(a)}]
Show that \(1-j\) is a unit for every \(j\in J\).
\item[\textnormal{(b)}]
Show that if~\(A\) is a finitely generated~\(R\)-module such that \(JA=A\), then \(A=0\). (Hint: Consider a generating set for~\(A\) of minimal size and derive a contradiction.)
\end{examsubparts}
\def\ProblemHeadingText{Problem 10.}
\item
Let~\(R\) be a ring (possibly non-commutative, and possibly without identity). Let~\(M\) be an irreducible left~\(R\)-module. Let \(D=\operatorname{End}_R(M)\).
\begin{examsubparts}
\item[\textnormal{(a)}]
Prove \(D=\operatorname{End}_R(M)\) is a division ring. (You may NOT assume Schur’s Lemma for this.)
\item[\textnormal{(b)}]
Give an example of a ring~\(R\) and an irreducible module~\(M\) such that~\(D\) is not commutative.
\end{examsubparts}
\def\ProblemHeadingText{Problem 11.}
\item
Classify (up to ring isomorphism) all semisimple rings of order \(720\).
\end{examproblems}
\end{document}
